\section*{Глава 3. Разработка веб-приложения}
\addcontentsline{toc}{section}{Глава 3. Разработка веб-приложения}
\stepcounter{section}

\subsection{Разработка базовой структуры приложения и вёрстка шаблонов страниц}

В соответствии с задачей 3.1, спроектирована компонентная структура клиентской части на базе шаблонизатора \textbf{Jinja2} с наследованием шаблонов. Базовый шаблон \texttt{base.html} определяет общий макет: адаптивный заголовок с навигацией, блок для flash-сообщений, основную область и подвал. Медиа-запросы и классы CSS (сетка \texttt{grid}, карточки \texttt{card}, адаптивные таблицы) обеспечивают корректное отображение на экранах разной ширины.

\subsubsection{Разработка шаблона страницы преподавателя и администратора, написание логики связи ``клиент-сервер''}

В рамках задачи 3.1.1 созданы шаблоны \texttt{teacher\_sessions.html}, \texttt{admin\_groups.html} и \texttt{admin\_users.html}. Ролевая фильтрация пунктов меню реализована через проверку \texttt{current\_user.role} внутри Jinja2 (листинг~\ref{lst:base_nav}).

\begin{lstlisting}[caption={Фрагмент навигации с ролевым разделением (base.html)}, label={lst:base_nav}]
{% if current_user.is_authenticated %}
  {% if current_user.role in ['teacher', 'admin'] %}
    <a href="{{ url_for('main.teacher_sessions') }}">Очереди и дисциплины</a>
    <a href="{{ url_for('main.journal') }}">Журнал</a>
  {% endif %}
  {% if current_user.role == 'admin' %}
    <a href="{{ url_for('main.admin_groups') }}">Группы</a>
    <a href="{{ url_for('main.admin_users') }}">Пользователи</a>
  {% endif %}
{% endif %}
\end{lstlisting}

Страница преподавателя (листинг~\ref{lst:teacher_form}) содержит две формы создания: дисциплины и очереди. Скрытое поле \texttt{action} позволяет серверу различать тип операции. Логика связи ``клиент-сервер'' реализована синхронными POST-запросами; после выполнения формы сервер возвращает редирект для обновления списка.

\begin{lstlisting}[caption={Фрагмент teacher\_sessions.html — создание дисциплины}, label={lst:teacher_form}]
<form method="post" class="form">
  <input type="hidden" name="action" value="subject">
  <label>Название<input name="name" required></label>
  <label>Группа<select name="group_id">...</select></label>
  <label>Преподаватель<select name="teacher_id">...</select></label>
  <label>Количество лабораторных работ
    <input type="number" name="lab_count" min="1" max="30" value="10" required>
  </label>
  <label>Минут на каждую лабораторную
    <input type="number" name="minutes_per_lab" min="1" max="120" value="15" required>
  </label>
  <button class="btn" type="submit">Создать</button>
</form>
\end{lstlisting}

\begin{figure}[ht!]
    \centering
    \includegraphics[width=0.85\textwidth]{screen_teacher_sessions}
    \caption{Страница преподавателя: форма создания дисциплины и список очередей}
    \label{fig:screen_teacher}
\end{figure}

Административные шаблоны используют аналогичный подход; например, \texttt{admin\_users.html} через \texttt{<select>} позволяет назначить роль и группу. Все операции над пользователями и группами защищены декоратором \texttt{@role\_required("admin")} (см. листинг~\ref{lst:role_decorator}).

\subsubsection{Разработка шаблона страницы студента и написание логики очереди}

Страница студента представлена шаблоном \texttt{dashboard.html} (список доступных очередей) и \texttt{session\_detail.html} (детальное отображение очереди). Форма записи в очередь (листинг~\ref{lst:join_form}) содержит поля: номер лабораторной, количество работ, предпочтение по времени, предполагаемая длительность, комментарий. Данные отправляются на маршрут \texttt{/sessions/<id>/join} методом POST.

\begin{lstlisting}[caption={Форма записи в очередь (session\_detail.html)}, label={lst:join_form}]
<form method="post" action="{{ url_for('main.join_queue', session_id=session.id) }}">
  <label>Номер лабораторной<input type="number" name="lab_number" value="1" required></label>
  <label>Количество работ<input type="number" name="work_count" value="1" required></label>
  <label>Предпочтение времени
    <select name="time_preference">
      <option value="any">Без разницы</option>
      <option value="early">Раньше</option>
      <option value="later">Позже</option>
    </select>
  </label>
  <label>Примерное время (мин)
    <input type="number" name="estimated_minutes"
           value="{{ session.subject.minutes_per_lab }}" required>
  </label>
  <label>Комментарий<input type="text" name="note"></label>
  <button class="btn" type="submit">Встать в очередь</button>
</form>
\end{lstlisting}

Логика очереди реализована на сервере в модуле \texttt{queue\_service.py} и описана в подразделе 3.3.1. Для обновления очереди в реальном времени используется WebSocket (см. \texttt{realtime.py}), который отправляет клиентам актуальное состояние после каждого изменения (добавление, удаление, вызов студента). На клиентской стороне JavaScript подключается к сокету \texttt{/ws/session/<id>} и динамически обновляет таблицу.

\begin{figure}[ht!]
    \centering
    \includegraphics[width=0.85\textwidth]{screen_student_dashboard}
    \caption{Главная страница студента: список доступных сессий для записи в очередь}
    \label{fig:screen_dashboard}
\end{figure}

\subsection{Реализация модуля авторизации и разграничения прав доступа}

Согласно задаче 3.2, в системе реализована ролевая модель с использованием расширения \textbf{Flask-Login}. Пароли хешируются с помощью \texttt{werkzeug.security}. Для защиты маршрутов и шаблонов применяются декоратор \texttt{@role\_required} и атрибуты объекта \texttt{current\_user}.

\subsubsection{Реализация модуля авторизации для Преподавателя и Студента}

Маршрут \texttt{/register} (листинг~\ref{lst:register}) принимает ФИО, email, пароль, роль и пригласительные коды для преподавателя (\texttt{TEACHER\_INVITE\_CODE}) или администратора (\texttt{ADMIN\_INVITE\_CODE}). Для студентов код не требуется. После успешной регистрации пользователь автоматически входит в систему.

\begin{lstlisting}[caption={Регистрация пользователя (auth.py)}, label={lst:register}]
if role == "teacher" and invite_code != current_app.config["TEACHER_INVITE_CODE"]:
    flash("Неверный код преподавателя.", "danger")
    return redirect(url_for("auth.register"))
user = User(full_name=full_name, email=email, role=role, group_id=group_id)
user.set_password(password)
db.session.add(user)
db.session.commit()
login_user(user)
\end{lstlisting}

\begin{figure}[ht!]
    \centering
    \includegraphics[width=0.75\textwidth]{screen_register}
    \caption{Страница регистрации: выбор роли и ввод пригласительного кода}
    \label{fig:screen_register}
\end{figure}

\begin{figure}[ht!]
    \centering
    \includegraphics[width=0.75\textwidth]{screen_login}
    \caption{Страница входа в систему}
    \label{fig:screen_login}
\end{figure}

Вход выполняется через маршрут \texttt{/login} (листинг~\ref{lst:login}), который проверяет email и пароль, после чего вызывает \texttt{login\_user}. Выход — \texttt{/logout} с удалением сессии. Преподаватели и студенты после входа перенаправляются на \texttt{dashboard}.

\begin{lstlisting}[caption={Аутентификация пользователя (auth.py)}, label={lst:login}]
user = User.query.filter(func.lower(User.email) == email).first()
if user and user.check_password(password):
    login_user(user, remember=bool(request.form.get("remember")))
    return redirect(next_url or url_for("main.dashboard"))
\end{lstlisting}

\subsubsection{Реализация модуля авторизации для администратора, разграничения прав доступа для всех ролей}

Для администратора создан отдельный маршрут регистрации с проверкой кода \texttt{ADMIN\_INVITE\_CODE}. Разграничение доступа на уровне маршрутов обеспечивается декоратором \texttt{@role\_required} (листинг~\ref{lst:role_decorator}), который принимает список допустимых ролей.

\begin{lstlisting}[caption={Декоратор role\_required (security.py)}, label={lst:role_decorator}]
def role_required(*roles: str):
    def decorator(view):
        @wraps(view)
        def wrapped(*args, **kwargs):
            if not current_user.is_authenticated or current_user.role not in roles:
                abort(403)
            return view(*args, **kwargs)
        return wrapped
    return decorator
\end{lstlisting}

Пример использования: \texttt{@role\_required("teacher", "admin")} перед маршрутами управления очередями и журналом. Для администратора доступны дополнительные маршруты: \texttt{/admin/groups}, \texttt{/admin/users}. В шаблонах проверка роли осуществляется через \texttt{current\_user.role} для условного отображения элементов интерфейса.

\subsection{Разработка CRUD интерфейса}

Реализованы полноценные CRUD-операции над основными сущностями: очереди, группы, пользователи, дисциплины. Взаимодействие клиент-сервер построено на HTTP-методах POST (создание, обновление) и GET (чтение), а также POST-запросах на удаление/изменение статуса.

\subsubsection{Разработка CRUD взаимодействия ``Клиент-Сервер'' для очередей}

Управление очередями осуществляется через маршруты, описанные в \texttt{main.py}. Основные операции:

\begin{itemize}
    \item \textbf{Создание очереди}: POST \texttt{/teacher/sessions} (с параметром \texttt{action=session}) — добавляет запись в таблицу \texttt{QueueSession}.
    \item \textbf{Запись студента}: POST \texttt{/sessions/<id>/join} — создаёт запись \texttt{QueueEntry}, вычисляет позицию через функцию \texttt{next\_position} (листинг~\ref{lst:next_pos}).
    \item \textbf{Вызов следующего}: POST \texttt{/sessions/<id>/call-next} — находит первый элемент в очереди со статусом \texttt{waiting} и переводит его в \texttt{checking} (функция \texttt{call\_entry}).
    \item \textbf{Перемещение записи}: POST \texttt{/entries/<id>/move/<direction>} — меняет позиции двух соседних записей.
    \item \textbf{Завершение проверки}: POST \texttt{/entries/<id>/pass} — сохраняет оценку в \texttt{GradeRecord}, меняет статус на \texttt{passed} и пересчитывает позиции.
\end{itemize}

\begin{figure}[ht!]
    \centering
    \includegraphics[width=0.85\textwidth]{screen_queue_detail}
    \caption{Страница активной очереди: список участников с кнопками управления}
    \label{fig:screen_queue}
\end{figure}

Логика очереди реализована в \texttt{queue\_service.py}. Для гарантии FIFO-порядка используется поле \texttt{position}, которое пересчитывается функцией \texttt{normalize\_positions} (листинг~\ref{lst:normalize}).

\begin{lstlisting}[caption={Пересчёт позиций в очереди}, label={lst:normalize}]
def normalize_positions(session_id: int) -> None:
    entries = QueueEntry.query.filter(
        QueueEntry.session_id == session_id,
        QueueEntry.status.in_(["waiting", "checking"])
    ).order_by(QueueEntry.position.asc(),
               QueueEntry.created_at.asc()).all()
    for index, entry in enumerate(entries, start=1):
        entry.position = index
\end{lstlisting}

Функция \texttt{next\_position} (листинг~\ref{lst:next_pos}) вычисляет максимальное значение позиции для активных записей и возвращает его, увеличенное на единицу.

\begin{lstlisting}[caption={Вычисление следующей позиции}, label={lst:next_pos}]
def next_position(session_id: int) -> int:
    max_pos = db.session.query(func.max(QueueEntry.position))
        .filter(QueueEntry.session_id == session_id,
                QueueEntry.status.in_(VISIBLE_QUEUE_STATUSES))
        .scalar()
    return int(max_pos or 0) + 1
\end{lstlisting}

\subsubsection{Разработка CRUD взаимодействия ``Клиент-Сервер'' для сущностей}

Для групп, пользователей и дисциплин реализованы стандартные CRUD-операции:

\begin{itemize}
    \item \textbf{Группы}: POST \texttt{/admin/groups} создаёт новую группу; GET отображает список всех групп.
    \item \textbf{Пользователи}: POST \texttt{/admin/users} создаёт пользователя с указанием роли, группы и пароля; GET — список всех пользователей.
    \item \textbf{Дисциплины}: POST \texttt{/teacher/sessions} (\texttt{action=subject}) создаёт дисциплину, привязывая её к группе и преподавателю.
\end{itemize}

Пример создания пользователя через административный интерфейс приведён в листинге~\ref{lst:admin_create_user}.

\begin{lstlisting}[caption={Создание пользователя администратором (main.py)}, label={lst:admin_create_user}]
user = User(email=email, full_name=full_name, role=role, group_id=group_id)
user.set_password(password)
db.session.add(user)
db.session.commit()
\end{lstlisting}

Все CRUD-операции защищены декораторами \texttt{@login\_required} и \texttt{@role\_required}, а также проверяют корректность входных данных перед записью в базу.

\subsection{Загрузка работ и выгрузка отчётов}

В соответствии с задачей 3.4, реализован функционал загрузки студенческих работ и выгрузки отчётов в формате CSV.

\textbf{Загрузка работ:} студент на странице профиля (\texttt{/profile}) через форму с \texttt{enctype="multipart/form-data"} отправляет файл, номер лабораторной и дисциплину. На сервере (листинг~\ref{lst:upload}) файл сохраняется в папке \texttt{UPLOAD\_FOLDER} с безопасным именем, а запись добавляется в таблицу \texttt{WorkUpload}. Скачать загруженный файл можно по маршруту \texttt{/uploads/<filename>} с проверкой прав доступа.

\begin{lstlisting}[caption={Обработка загрузки работы (main.py)}, label={lst:upload}]
original = secure_filename(file.filename)
safe_name = f"{current_user.id}_{uuid4().hex}_{original}"
upload_path = Path(current_app.config["UPLOAD_FOLDER"]) / safe_name
file.save(upload_path)
upload = WorkUpload(
    student_id=current_user.id,
    subject_id=subject_id,
    lab_number=lab_number,
    filename=safe_name,
    original_filename=original,
    comment=comment
)
db.session.add(upload)
\end{lstlisting}

\begin{figure}[ht!]
    \centering
    \includegraphics[width=0.85\textwidth]{screen_profile_upload}
    \caption{Страница профиля студента: форма загрузки файла лабораторной работы}
    \label{fig:screen_upload}
\end{figure}

\textbf{Выгрузка отчётов:} для преподавателя и администратора доступны два CSV-отчёта:

\begin{itemize}
    \item \texttt{/reports/journal.csv} — журнал успеваемости: студент, группа, предмет, лабораторная, оценка, дата, комментарий.
    \item \texttt{/reports/queue/<session\_id>.csv} — состояние очереди: позиция, студент, номер ЛР, количество работ, предпочтение, время, подтверждение, статус.
\end{itemize}

Генерация CSV выполняется через модуль \texttt{csv} и возвращается как \texttt{Response} с типом \texttt{text/csv}. Пример формирования отчёта очереди приведён в листинге~\ref{lst:csv_report}.

\begin{lstlisting}[caption={Генерация CSV-отчёта очереди}, label={lst:csv_report}]
output = StringIO(newline="")
output.write("sep=;\r\n")  # директива Excel: разделитель — точка с запятой
writer = csv.writer(output, delimiter=";", lineterminator="\r\n")
writer.writerow(["Позиция", "Студент", "Лабораторная", "Кол-во работ",
                 "Предпочтение", "Время", "Подтвержден", "Статус"])
for entry in session.entries:
    writer.writerow([entry.position, entry.student.full_name, entry.lab_number,
                     entry.work_count, entry.preference_title, entry.estimated_minutes,
                     "да" if entry.is_confirmed else "нет", entry.status_title])
csv_bytes = ("" + output.getvalue()).encode("utf-8")  # BOM для Excel
return Response(csv_bytes, mimetype="text/csv; charset=utf-8",
                headers={"Content-Disposition": f"attachment; filename=queue_{session.id}.csv"})
\end{lstlisting}

\subsection{Вёрстка адаптивного интерфейса (Jinja2, HTML, CSS)}

Финальный этап разработки (задача 3.5) заключался в создании адаптивного пользовательского интерфейса с использованием \textbf{HTML5}, \textbf{CSS3} и шаблонизатора \textbf{Jinja2}. Основные принципы:

\begin{itemize}
    \item Использование CSS-классов для каркаса: \texttt{.container} (центрирование, максимальная ширина), \texttt{.grid} (адаптивная сетка), \texttt{.card} (карточка с тенью и скруглениями).
    \item Адаптивные таблицы: обёртка \texttt{.table-wrap} с горизонтальной прокруткой на узких экранах.
    \item Мобильное меню: кнопка \texttt{.nav-toggle} (JavaScript переключает класс \texttt{.nav--open} для отображения навигации).
    \item Формы используют классы \texttt{.form}, \texttt{.inline-form}, \texttt{.compact} для компактного отображения на мобильных устройствах.
    \item Цветовая схема и отступы обеспечивают читаемость на экранах от 320px до 1920px.
\end{itemize}

Пример CSS-класса для карточек очереди (листинг~\ref{lst:css_card}) демонстрирует адаптивные свойства: \texttt{flex} или \texttt{grid}, скругление, тень.

\begin{lstlisting}[caption={CSS-стили для карточек и адаптивных таблиц}, label={lst:css_card}]
.card {
    background: var(--surface);
    border: 1px solid var(--border);
    border-radius: 24px;
    box-shadow: var(--shadow);
    padding: 24px;
}
@media (max-width: 820px) {
    .hero, .grid.two, .grid.three { grid-template-columns: 1fr; }
    /* Адаптивные таблицы: карточки вместо горизонтальной прокрутки */
    .table-wrap { overflow-x: visible; border: none; }
    .table-wrap thead { display: none; }
    .table-wrap tr {
        display: block;
        border: 1px solid var(--border);
        border-radius: 16px;
        margin-bottom: 12px;
    }
    .table-wrap td[data-label]::before {
        content: attr(data-label);
        font-weight: 700;
        font-size: .78rem;
        text-transform: uppercase;
    }
}
\end{lstlisting}

\begin{figure}[ht!]
    \centering
    \includegraphics[width=0.85\textwidth]{screen_queue_desktop}
    \caption{Страница очереди на широком экране: таблица с полным набором столбцов}
    \label{fig:screen_desktop}
\end{figure}

\begin{figure}[ht!]
    \centering
    \includegraphics[width=0.5\textwidth, height=0.65\textheight, keepaspectratio]{screen_queue_mobile}
    \caption{Страница очереди на мобильном устройстве (ширина экрана $\leq$820px): карточный вид}
    \label{fig:screen_mobile}
\end{figure}

Навигационное меню в \texttt{base.html} содержит кнопку-гамбургер и JavaScript-обработчик для адаптивного отображения на мобильных устройствах. Все интерактивные элементы (кнопки, формы) имеют достаточный размер для удобного нажатия пальцем.

Таким образом, разработанное веб-приложение полностью соответствует требованиям технического задания: реализованы ролевая модель, FIFO-очередь с подтверждением, CRUD-интерфейс, загрузка файлов, выгрузка отчётов и адаптивная вёрстка на Jinja2/HTML/CSS. Исходный код развёртывается через Docker Compose или локально с PostgreSQL.


\subsection{Реализация WebSocket-обновления очереди в реальном времени}

Для обеспечения синхронного отображения очереди у всех подключённых пользователей (студентов, преподавателя, администратора) реализован механизм реального времени на базе протокола \textbf{WebSocket} с использованием библиотеки \textit{Flask-Sock}~\cite{flask_docs}. При изменении состояния очереди (добавление, удаление, перемещение, вызов студента) сервер инициирует широковещательную рассылку актуальных данных всем клиентам, подписанным на данную сессию.

Схема работы представлена на листинге~\ref{lst:websocket_broadcast}. Функция \texttt{broadcast\_session\_update} сериализует текущее состояние очереди с помощью \texttt{serialize\_session} (определена в \texttt{queue\_service.py}) и отправляет его во все клиентские очереди, зарегистрированные для указанного \texttt{session\_id}.

\begin{lstlisting}[caption={Широковещательная рассылка обновлений (realtime.py)}, label={lst:websocket_broadcast}]
_clients: dict[int, set[Queue]] = defaultdict(set)
_lock = RLock()

def broadcast(session_id: int, payload: dict[str, Any]) -> None:
    with _lock:
        targets = list(_clients.get(session_id, set()))
    for client_queue in targets:
        client_queue.put(payload)

def broadcast_session_update(session_id: int) -> None:
    session = QueueSession.query.get(session_id)
    if session:
        broadcast(session_id, serialize_session(session))
\end{lstlisting}

WebSocket-маршрут \texttt{/ws/session/<int:session\_id>} (листинг~\ref{lst:ws_route}) принимает подключение от клиента, регистрирует его очередь сообщений в глобальном словаре и отправляет начальный снимок состояния. При каждом обновлении сервер помещает сериализованный JSON в очередь клиента, а клиентский JavaScript получает сообщение и динамически обновляет таблицу очереди без перезагрузки страницы.

\begin{lstlisting}[caption={Обработчик WebSocket-соединения (realtime.py)}, label={lst:ws_route}]
@sock.route("/ws/session/<int:session_id>")
def session_socket(ws, session_id: int):
    if not current_user.is_authenticated:
        ws.close()
        return
    client_queue: Queue = Queue()
    with _lock:
        _clients[session_id].add(client_queue)
    try:
        session = QueueSession.query.get(session_id)
        if session:
            ws.send(json.dumps(serialize_session(session)))
        while True:
            payload = client_queue.get(timeout=25)
            ws.send(json.dumps(payload, ensure_ascii=False, default=str))
    except Exception:
        pass
    finally:
        with _lock:
            _clients[session_id].discard(client_queue)
\end{lstlisting}

На клиентской стороне JavaScript (файл \texttt{static/js/main.js}) при загрузке страницы очереди устанавливает WebSocket-соединение. При получении сообщения типа \texttt{"snapshot"} таблица очереди перерисовывается. Для браузеров, не поддерживающих WebSocket, или в случае разрыва соединения предусмотрен резервный механизм периодического опроса через API-маршрут \texttt{/api/sessions/<id>/snapshot} (листинг~\ref{lst:api_snapshot}).

\begin{lstlisting}[caption={REST-совместимый снимок очереди (api.py)}, label={lst:api_snapshot}]
@api_bp.route("/sessions/<int:session_id>/snapshot")
@login_required
def session_snapshot(session_id: int):
    session = QueueSession.query.get_or_404(session_id)
    return jsonify(serialize_session(session))
\end{lstlisting}

Таким образом, достигнута синхронизация состояния очереди в реальном времени, что критически важно для удобства работы преподавателя и студентов в ходе занятия.